\cotefile{14}
\scanpage{16}
\mspage{II. 9}

\[
\handwritten{\overline{\zeta\eta}}: \underline{\mathrm{Hom}}_{\handwritten{\mathcal{A}}}(\handwritten{\mathcal{M}},\handwritten{\unsure{\mathcal{M}|\mathcal{M}\mathcal{O}}})\otimes_{\handwritten{\mathcal{A}}}\underline{\mathrm{Hom}}_{\handwritten{\mathcal{A}}}(\handwritten{\mathcal{N}},\handwritten{\mathcal{N}})\longrightarrow \underline{\mathrm{Hom}}_{\handwritten{\mathcal{A}}}(\handwritten{\mathcal{M}}\otimes_{\handwritten{\mathcal{A}}}\handwritten{\mathcal{N}},\handwritten{\mathcal{M}}\otimes_{\handwritten{\mathcal{A}}}\handwritten{\mathcal{N}}).
\]

Since $\handwritten{\mathcal{M}}$ and $\handwritten{\mathcal{N}}$ are finitely presented, given any $x$ in $X$, $\handwritten{\overline{\zeta\eta}}_x$ is
a homomorphism from

\[
\begin{aligned}
&\underline{\mathrm{Hom}}_{\handwritten{\mathcal{A}}_x}(\handwritten{\mathcal{M}}_x,\handwritten{\mathcal{M}}_x)\handwritten{\otimes}_{\handwritten{\unsure{\mathcal A_x}}}\underline{\mathrm{Hom}}_{\handwritten{\mathcal{A}}_x}(\handwritten{\mathcal{N}}_x,\handwritten{\mathcal{N}}_x)\quad\text{to}\\
&\underline{\mathrm{Hom}}_{\handwritten{\mathcal{A}}_x}(\handwritten{\mathcal{M}}_x\handwritten{\otimes}_{\handwritten{\mathcal A_x}}\handwritten{\mathcal{N}}_x,\handwritten{\mathcal{M}}_x\handwritten{\otimes}_{\handwritten{\mathcal A_x}}\handwritten{\mathcal{N}}_x).
\end{aligned}
\]

Since $\handwritten{\mathcal{M}}_x$ and $\handwritten{\mathcal{N}}_x$ are finitely generated projective $\handwritten{\mathcal{A}}_x$-modules,
$\handwritten{\overline{\zeta\eta}}_x$ is an isomorphism for every $x$ in $X$. Thus $\handwritten{\overline{\zeta\eta}}$ is an isomorphism of sheaves.

If $(\handwritten{\mathcal{M}},\handwritten{\mu})$ is a left module over $(\handwritten{\mathcal{A}},\alpha)$ we can consider
$\underline{\mathrm{Hom}}_{\handwritten{\mathcal{A}}}(\handwritten{\mathcal{M}},\handwritten{\mathcal{A}})$ a right $\handwritten{\mathcal{A}}$-module by means fo the operation of
$\handwritten{\mathcal{A}}(U)$ on $\mathrm{Hom}_{\handwritten{\mathcal{A}}|U}(\handwritten{\mathcal{M}}|U,\handwritten{\mathcal{A}}|U)$ defined for any open set $U$ of $X$
as follows : for $f=\{f_V\}_{V\in U}$ and $a$ in $\handwritten{\mathcal{A}}(U)$, $f\cdot a=g=\{g_V\}_{V\subset U}$
where $g_V(m)=f_V(m)\,\alpha^U_V(a)$ for any $m$ in $\handwritten{\mathcal{M}}(V)$.

Now we can define a mapping

\[
\psi:\underline{\mathrm{Hom}}_{\handwritten{\mathcal{A}}}(\handwritten{\mathcal{M}},\handwritten{\mathcal{A}})\handwritten{\otimes}_{\handwritten{\mathcal{A}}}\handwritten{\mathcal{M}}\longrightarrow \underline{\mathrm{Hom}}_{\handwritten{\mathcal{A}}}(\handwritten{\mathcal{M}},\handwritten{\mathcal{M}}).
\]

First we define the homomorphism of presheaves
$\psi':\underline{\mathrm{Hom}}_{\handwritten{\mathcal{A}}}(\handwritten{\mathcal{M}},\handwritten{\mathcal{A}})\handwritten{\otimes}_{\handwritten{\mathcal{A}}}^{\handwritten{\prime}}\handwritten{\mathcal{M}}\longrightarrow \underline{\mathrm{Hom}}_{\handwritten{\mathcal{A}}}(\handwritten{\mathcal{M}},\handwritten{\mathcal{M}})$ in the following way : for
$f=\{f_V\}_{V\subset U}$ in $\mathrm{Hom}_{\handwritten{\mathcal{A}}|U}(\handwritten{\mathcal{M}}|U,\handwritten{\mathcal{A}}|U)$ and $m$ in $\handwritten{\mathcal{M}}(U)$,

% archive sheet number: 15
\scanpage{17}
\unsure{Équivalences|Equivalences} d'Albanese et de Picard des cycles algébriques

% archive sheet number: 16
\scanpage{19}
\deleted{$E$}

$A$ courbe elliptique

\[
Y=A\times P^{\deleted{\unsure{k|2|1}}1}\qquad\text{$P$ courbe non. sing. complète conn.}
\]

\[
\begin{aligned}
\Theta_0={}&\Delta-A\times(0)-(0)\times A\quad\text{\deleted{diviseur sur} diviseur can. sur } A\times A\\
&\text{définissant l'autodualité}
\end{aligned}
\]

\[
\Theta=\Theta_0\times P^{\deleted{\unsure{k|1}}1}_{p}\qquad\text{diviseur correspondant sur }A\times A\times P^{\unsure{1|k}}=A\times Y
\]

$c$, classe de diviseur sur $Y$, défini par un \deleted{\unsure{fibré}} faisceau inv. $\underline L$ comme $c=\mathrm{cl}(\underline L)$

\[
X=\overline{V(\underline L)}\quad\text{(clôture projective)}.\quad Y^2\xrightarrow{i}X^3\ \text{section nulle.}
\]

le cycle $\Theta$ sur $A\times Y$ définit

\[
z=(\mathrm{id}_A\times i')_*(\Theta)
\]

cycle de codim $2$ sur $\deleted{\unsure{X}}\,A\times X^3$. On veut calculer, si $s:A\times X^3\unsure{\longrightarrow|-}A$ est la projection, $s_*(z^2)$, et on va donner (pour \deleted{des} choix convenables de $c$) des cas où il est $>0$, d'autres où il est $<0$.

Pour ceci, notons que si $i'=\mathrm{id}_A\times i$, on a

\diagram{\begin{tikzcd}[ampersand replacement=\&,column sep=small]
A\times Y \arrow[dr,"t"] \& \& A\times X \arrow[ll,"i'"'] \arrow[dl,"s"]\\
\& A \&
\end{tikzcd}}

\[
\begin{aligned}
z^2&=i'_*(\Theta)^2=i'_*(\Theta^2c')\\
\text{où}\quad c'&=i'^{*}(i'_*(1))=c^1\ (\text{faisceau normal pour }i')\\
&=1_A\times c,\quad c\ \text{étant le faisceau normal}\\
&\hspace{7em}\text{\unsure{pour|par} }\deleted{\unsure{Y}}\,i:Y\unsure{-|\longrightarrow} X.
\end{aligned}
\]

% archive sheet number: 18
\scanpage{20}
\mspage{- 15 -}

n\textsuperscript{o} 391

d'où $s'(x)=s(x)+2t$, et en remplaçant $x$ par $s(x)$, il vient $s'(s(x))=x+2t$.
Dans le cas particulier où $K=R$ et où $\Phi$ est positive non dégénérée, on vérifie aussitôt que $\lVert t\rVert$ est la \uline{distance} des hyperplans $H$, $H'$ (pour la structure d'espace métrique définie sur $E$ par la distance euclidienne).

5. \uline{Chambres dans un espace affine}.

Soient $E$ un \uline{espace affine réel} de dimension finie $n$ ; en prenant une origine dans $E$, $E$ s'identifie à un espace vectoriel réel de dimension $n$ et est donc muni d'une topologie bien déterminée pour lequel il est isomorphe à $R^n$ (\uline{Top}. \uline{gén}., chap. VI, §1, n\textsuperscript{o} 5) ; en outre, les translations étant des homéomorphismes de $R^n$ sur lui-même, cette topologie est indépendante de l'origine choisie dans $E$ ; quand nous considérerons dans ce qui suit $E$ comme espace topologique, il s'agira toujours de cette topologie.

Pour tout hyperplan $H$ dans $E$ et tout ensemble \uline{connexe} non vide $C$ ne rencontrant pas $H$, $C$ est contenu dans un des demi-espaces ouverts déterminés par $H$, ces derniers formant une partition de $E-H$ en deux ensembles ouverts connexes ; on désignera le demi-espace ouvert déterminé par $H$ et contenant $C$ par $D_H(C)$.

LEMME 1.- \uline{Si} $\underline H$ \uline{est un ensemble localement fini} (\uline{Top}. \uline{gén}., chap. I, 3\textsuperscript{e} éd., § 1, n\textsuperscript{o} 5, déf. 8) \uline{d'hyperplans affines de} $E$, \uline{la réunion} $\unsure{S|S'}$ \uline{des} $H\in\underline H$ \uline{est un ensemble fermé et rare dans} $E$, \uline{et pour toute composante connexe} $C$ \uline{de} $E-S$, \uline{on a} $C=\bigcap\limits_{H\in\underline H}D_H(C)$.

Comme les $H\in\underline H$ sont fermés, on sait qu'il en est de même de $S$ (\uline{Top}. \uline{gén}., chap. I, 3\textsuperscript{e} éd., §1, n\textsuperscript{o} 5, prop. 4) ; d'autre part, chacun des $H\in\underline H$ est rare dans $E$ ; si $x\in E$, il y a un voisinage ouvert $V$ de $x$ ne rencontrant qu'un nombre fini d'éléments $H_i\in\underline H$ $(1\leq i\leq n)$, et la réunion des $V\cap H_i$ est rare dans $V$, puisqu'il en est ainsi de chacun d'eux. Pour tout $H\in\underline H$, il est clair

% archive sheet number: 19
\scanpage{21}
D'où $s_*(z^2)=t_*(\Theta^2c')$, où $t:A\times Y\unsure{-|\longrightarrow}A$ est la projection. On a alors, du pt de vue \unsure{degré}

\diagram{\begin{tikzcd}[ampersand replacement=\&] A\times Y\arrow[d,"t"']\\ A\end{tikzcd}}

\[
\begin{aligned}
\unsure{\mathrm{tr}(z^2)|\mathrm{tr}(s_*(z^2))|K(z^2)}&=\mathrm{tr}(\Theta^2c')=\mathrm{tr}(\Theta^2\cdot1_A\times c)\\
&=\mathrm{tr}(p_*(\Theta^2)\cdot c)
\end{aligned}
\]

où $p$ est la projection $A\times Y\unsure{-|\longrightarrow}Y$.

Nous allons calculer $p_*(\Theta^{\unsure{2|?}})$, en décomposant $p$ en $\deleted{\unsure{q}}p'\times\mathrm{id}_{P}$, où $p':A\times A\longrightarrow A$ est la deuxième projection. On a

\diagram{\begin{tikzcd}[ampersand replacement=\&,row sep=small]
 A\times\deleted{\unsure{X}}Y\arrow[r,"p"] \arrow[d,equal] \&Y\\
 A\times\underbrace{A\times P}\arrow[r] \&\underbrace{A\times P}
\end{tikzcd}}

\[
\begin{aligned}
\Theta&=\Theta_0\times1_P\quad\text{d'où }\Theta^2=\Theta_0^2\times1_P\qquad\text{d'où}\\
p_*(\Theta^2)&=p'_*(\Theta_0^2)\times\mathrm{id}_{P*}(1_P)=p'_*(\Theta_0^2)\times1_P.
\end{aligned}
\]

D'autre part

\[
\begin{aligned}
\Theta_0^2&=\deleted{\Delta}\,\Delta^2+0+0-2(0,0)-2(0,0)+2(0,0)\\
&=-2(0,0)
\end{aligned}
\]

compte tenu que $\Delta^2=0$ \inserted{(équiv. alg.)} puisque $\chi(A)=0$.

d'où

\[
p_*(\Theta^2)=-2(0\times0\times P)
\]

et on trouve par suite

\[
\mathrm{tr}(s_*(z^2))=-2c(0,0)
\]

où on pose

\[
c(0,0)=\mathrm{tr}\bigl(c\cdot(0\times\unsure{0|\varnothing}\times P)\bigr)
\]

Il existe un $c$ tel que $c(0,0)=1$, p. ex. $c=A\times A\times(a)$, où $a\in P(k)$. Donc remplaçant \inserted{\unsure{ensuite}} $c$ par $-c$, on trouve comme valeurs successivement $-2$ et $2$, qui ont des signes différents !!

% archive sheet number: 20
\scanpage{22}
\mspage{- 14 -}

n\textsuperscript{o} 391

\[
(t,\lambda)\handwritten{\longrightarrow}({}^{v}u(t)+\lambda(u(a)-a'),\lambda).\tag*{(9)}
\]

Il est clair que l'application $u\handwritten{\longrightarrow}{}^{h}u$ est \uline{injective}.

On appelle \uline{pseudo-réflexion affine} une application affine $s$ de $E$ dans lui-même qui laisse fixes les points d'un \uline{hyperplan \unsure{affine|af'ine}} $H$. Si on prend pour origine dans $E$ un point $a\in H$, $s$ devient une \uline{pseudo-réflexion} au sens du n\textsuperscript{o} 2 ; donc $^{v}s$ est une pseudo-réflexion laissant invariants les vecteurs de l'hyperplan $^{v}H$, \uline{direction} de $H$ ; et la formule (9) montre que $^{h}s$ est une \uline{pseudo-réflexion} dans l'espace vectoriel $^{h}E$, laissant invariants les points de l'hyperplan $^{h}H$ \uline{engendré par} $H$.

On dit qu'une pseudo-réflexion affine $s$ est une \uline{réflexion affine} si elle est involutive et distincte de \deleted{l'}identité ; il revient au même de dire que la pseudo-réflexion $^{v}s$ ou $^{h}s$ est une \uline{réflexion} (n\textsuperscript{o} 2), comme il résulte aussitôt de ce qui précède.

Lorsque $K$ est de caractéristique $\neq2$ et que $T$ est muni d'une forme bilinéaire symétrique non dégénérée $\Phi$, les réflexions de $E$ qui sont des \uline{déplacements} pour la forme métrique $\Phi$ (\uline{Alg}., chap. IX, §6, n\textsuperscript{o} 6) sont les \uline{symétries orthogonales par rapport aux hyperplans affines} (non isotropes) de $E$ (n\textsuperscript{o} 2, prop. 4).

\margin{\handwritten{\unsure{/}}}

\margin{\handwritten{\unsure{!}}}

PROPOSITION 9.- \uline{Soient} $H$, $H'$ \uline{deux hyperplans dans} $E$ \uline{parallèles et non isotropes} (\uline{pour} $\Phi$) ; \uline{il existe un unique vecteur} $t\in T$ \uline{orthogonal à} $H$ \uline{et tel que} $H'=t+H$. \uline{Si} $s$, $s'$ \uline{sont les symétries orthogonales par rapport à} $H$ \uline{et} $H'$ \uline{respectivement}, $s's$ \uline{est la translation} $2t$ \uline{dans} $E$.

En prenant dans $E$ une origine contenue dans $H$, on peut se borner au cas où $E=T$. L'existence et l'unicité de $t$ sont immédiates. En prenant une nouvelle origine au point $t$, on voit que

\[
s'(x)-t=s(x-t)=s(x)+t
\]

% archive sheet number: 21
\scanpage{24}
\margin{Consulter Weil, \unsure{Critères} d'\unsure{Équivalence|Équivalences}}

\uline{Théorie de l'équivalence d'Albanese des cycles}

1. \uline{L'anneau \deleted{\unsure{de l'|des}}des \inserted{classes d'}\unsure{équivalence|équivalences} d'Albanese}

\uline{Définition} 1.1. Soit $X$ un schéma \deleted{[projectif]} lisse sur le corps $k$, $z$ un cycle sur $X$. On dit que $z$ est alb-équiv. à $0$,

\margin{?}

si on peut trouver un schéma [projectif] lisse \deleted{\unsure{irréductible}}\inserted{connexe} $\deleted{\unsure{Y}}T$ sur $k$, \deleted{\unsure{et}}\deleted{\unsure{et}} un $0$-cycle \deleted{de degré $0$ $\unsure{\mathfrak a|\alpha|\mathfrak{t}}$} $\mathfrak a$ sur $\deleted{\unsure{Y}}T$, et un cycle $Z$ sur $X\times T$, tels que

(i) \inserted{$\mathfrak a$ est de degré $0$ et} $S_Y(\mathfrak a)=0$\qquad (ii) $Z(\mathfrak a)$ est défini et égal à $z$.

\uline{Proposition} 1.2. (i) Les cycles alb-équiv. à $0$ sur $X$ forment un sous-groupe \inserted{gradué $N^*(X)$} de $\mathrm{Cyc}^*(X)$, contenant le groupe des cycles lin. équiv. à $0$, de sorte que la connaissance de $N^*(X)$ équivaut à celle de son image $N'{}^*(X)$ dans \deleted{Cyc} $A^*(X)$ (l'anneau de Chow de $X$, défini par les classes à équiv. linéaire près), et $\mathrm{Cyc}(X)/N(X)\simeq\deleted{\unsure{\mathrm{Cyc}}}A^*(X)/N'{}^*(X)$.

% archive sheet number: 23
\scanpage{25}
(ii) $N'{}^*(X)$ est un \uline{idéal} de $A^*(X)$, donc $\mathrm{Cyc}^*(X)/N^*(X)\simeq A^*(X)/N'{}^*(X)$ est muni d'une structure d'\uline{anneau}, soit $A'{}^*(X)$

(iii) Pour tout hom $f:X\longrightarrow Y$ de \deleted{variétés projectives}\inserted{\unsure{qu. pr.|qu. proj.|q.p.} \deleted{\unsure{préschémas}}} lisses, l'hom $f^*:A^*(Y)-A^*(X)$ applique $N'{}^*(Y)$ dans $N'{}^*(X)$, donc définit par passage au quotient un hom d'anneaux $A'{}^*(Y)-A'{}^*(X)$.

(iv) Pour tout $f$ comme dans (iii), \inserted{avec $f$ \uline{propre},} l'hom \inserted{de groupes} $f_*:A^*(X)\unsure{\longrightarrow|-}A^*(Y)$ \deleted{est} applique $N'{}^*(X)$ dans $N'{}^*(Y)$, et induit donc un hom de groupes (également noté $f_*$)

\[
A'{}^*(X)\longrightarrow A'{}^*(Y)
\]

(v) L'équiv. d'Albanese $\Longrightarrow$ équivalence algébrique. (Donc on peut parler d'\unsure{éléments|élémts} de $A'{}^*(X)$ \uline{alg.} \uline{équiv.} \uline{à} $0$ !)

\uline{Dém}. (i) \deleted{Si $z$ est alb-équiv. \unsure{à} $0$ sur $X$, il est évident qu'il en est de $\unsure{\overline{\mathrm m}|=}$ de $-z$.} \inserted{\unsure{Évident|Évid\textsuperscript{mt}|Évidemment} $0$ est alb-équiv. à $0$.}

Prouvons que si $z,z'$ sont alb-équiv. à $0$ sur $X$, il en est de $\overline{\mathrm m}$ de $z\unsure{-|+}z'$.

Soient $T,T'$ des variétés de paramètres comme dans 1.1. correspondant à $z,z'$. Soit $T_{1}^{\unsure{\prime\prime}}=T\times T'$

% archive sheet number: 24
\scanpage{26}
Alors $Z,Z'$ ont des images inverses $Z_1,Z'_1$ sur $X\times T_1$, et si $\unsure{\mathfrak a|\alpha|\mathfrak{t}},\mathfrak a'$ sont les $0$-cycles de $T,T'$ comme dans 1.1., on peut les \unsure{remonter} en $\mathfrak a_1=\mathfrak a\times a'$, $\mathfrak a'_1=a\times\mathfrak a'$ de $T\times T'$. Notons que

\margin{Lemme à dégager}

\[
S_{T\times T'}(\mathfrak a\times a')=S_Y(\mathfrak a)=0,\quad S_{T\times T'}(a\times\mathfrak a')=S_{T'}(\mathfrak a')=0
\]

donc $\mathfrak a_1^{\unsure{*|\prime}}$ et $\mathfrak a'_1$ sont \deleted{alb-équiv}\inserted{de degré $0$ et} tels que leurs $S_{T\times T'}$ sont nuls. D'autre part, on voit de suite que

\[
Z_1(\mathfrak a_1)=Z(\mathfrak a),\quad Z'(\mathfrak a'_1)=Z'(\mathfrak a'),
\]

\deleted{donc on est} d'autre part

\[
Z_1(\mathfrak a'_1)=Z(\mathrm{pr}_{1*}(\mathfrak a'_1))=0
\]

car $\mathrm{pr}_{1*}(\mathfrak a'_1)=0$, sous réserve que $\deleted{\unsure{Z}}a$ soit choisi tel que $Z(a)$ soit défini. On a de $\overline{\mathrm m}$, si $a'$ est choisi convenablement, $Z'_1(\mathfrak a_1)=0$.

\diagram{\begin{tikzcd}[ampersand replacement=\&,row sep=small,column sep=small]
X\times T \arrow[d,no head] \& X\times T\times T'\arrow[l]\arrow[d,no head]\\
T \& T\times T'\arrow[l]
\end{tikzcd}}

On posera alors

\[
Z_2=Z_1+Z'_1,\qquad\mathfrak a_2=\mathfrak a_1+\mathfrak a'_1
\]

et on aura

\[
Z_2(\mathfrak a_2)=z+z'.
\]

\deleted{\unsure{cqfd}} Donc $N(X)$ est bien un \unsure{ss-|s-}groupe de $X$.

% archive sheet number: 25
\scanpage{27}
Le fait qu'il soit gradué est immédiat. Il contient le groupe des cycles lin. équiv. à $0$, comme il résulte aussitôt \unsure{des définitions|de définition}, et du fait que sur $T=\mathbb P^1$, tout $0$-cycle $\unsure{\mathfrak a|\alpha|\mathfrak{t}}$ de degré $0$ satisfait $S_T(\mathfrak a)=0$. \deleted{Ainsi $N(X)$ est bien un \unsure{ss-|s-}groupe de $X$, contenant}

Pour prouver (iii), notons $z$ un cycle de la classe envisagée dans $A^*(\deleted{\unsure{T}}Y)$, et considérons un $Z$ sur $Y\times T$ et un $\mathfrak a$ comme dans 1.1., donc $z=Z(\mathfrak a)$. Si nous changeons $Z,\mathfrak a$ dans leurs classes d'équivalence linéaire \inserted{en $Z',\mathfrak a'$}, $Z'(\mathfrak a)$ sera \unsure{linéair\textsuperscript{mt}|linéairement} équivalent à $z$, et le Pb est le $\overline{\mathrm m}$ pour $z$ et $z'$.

Il suffit maintenant de prouver \deleted{$\unsure{Z'}$} qu'il existe $Z',\mathfrak a'$ tels que $(f\times\mathrm{id}_T)^*(Z')=Z'_1$ et $Z'(\mathfrak a'),\deleted{\unsure{Z_1(\mathfrak a')|Z'_1(\mathfrak{t}')}}$ soient définis ;

\diagram{\begin{tikzcd}[ampersand replacement=\&,column sep=small,row sep=small]
X\arrow[rr] \& \& Y\\
X\times T\arrow[rr,leftrightarrow]\arrow[dr] \& \&Y\times T\arrow[dl]\\
\deleted{\unsure{T}} \& T \&
\end{tikzcd}}

en effet, en vertu de la théorie des intersections, on aura alors $f^*(Z'(\mathfrak a'))=\deleted{\unsure{f^*}}(Z_1(\mathfrak a')$, or le premier mb \deleted{\unsure{représente est dans la|représentant et dans la}} représente $f^*(z)$ dans $A^*(X)$, \deleted{le deuxième est}

% archive sheet number: 26
\scanpage{28}
% slanted
\margin{à \uline{vérifier}}

Nous laissons ici le détail de la vérification de l'existence de $Z',\unsure{\mathfrak a'|\alpha'|\mathfrak{t}'}$.

On prouve (ii) par \deleted{\unsure{la}}\inserted{$\overline{\mathrm m}$} méthode que (ii), \deleted{\unsure{on en signera|à ce sujet}} en se ramenant au cas où $z,z'$ sont définis en termes des $\overline{\mathrm m}$ $T,\mathfrak a$ et deux cycles $Z,Z'$ sur $T$. On \inserted{peut} se ramener \unsure{aussi|ainsi} à (iii) et à l'énoncé analogue à (ii) pour les produits \uline{cartésiens}, ce qui dispense de faire encore bouger \unsure{des|les} cycles.

La démonstration de (iv) se fait par l'AQT. Celle de (v) est triviale sur les définitions (ou sur le formalisme de l'équiv. alg.)

\deleted{\uline{\unsure{Théorème}}} \uline{Scholie} 1.3. Les $A'{}^*(X)$, pour $X$ variable, ont les $\overline{\mathrm m}$s propriétés fonctorielles générales que les $A^*(X)$. En particulier, on a la formule de projection, Dans les cas où équiv. linéaire $=$ équiv. alg. sur $X$ (p. ex. l'espace ponctuel, les $\mathbb P^n$, les variétés de drapeaux) ils sont $=$ équiv. alb.

% archive sheet number: 27
\scanpage{29}
% slanted
\margin{Ici il semble qu'on utilise le fait que dans 1.1. on prend $T$ \uline{projectif}}

\uline{Théorème} 1.4. Soit $X$ projective lisse connexe de dim $n$ sur $k$. \deleted{Alors pour les cycles de dim $0$ et pour les cycles de codim $1$, l'équiv.}

(i) Soit $\unsure{\mathfrak a|\alpha|\mathfrak{t}}$ un $0$-cycle sur $X$. Alors $\mathfrak a$ est alb. équiv. à $0$ \unsure{sss|ssi} $\mathfrak a$ est de degré $0$, et $S_{\deleted{\unsure{T}}X}(\mathfrak a)=0$.

(ii) Soit $D$ un diviseur sur $X$. Alors $D$ est alb-équiv. à $0$ ssi $D$ est lin. équivalent à $0$.

\uline{Remarque} 1.5. On ignore, \inserted{déjà} pour les $0$-cycles, si l'équiv. d'Alb. est identique à l'équivalence linéaire (vrai pour tous les cycles ???).

\uline{Démonstration} (i) Si $\mathfrak a$ est de degré $0$ et $S_X(\mathfrak a)=0$, on a $\mathfrak a$ alb-équiv. à $0$ par définition [prendre $T=X$, $Z=$ diagonale de $X$, donc $\mathfrak a=Z(\mathfrak a)$]. L'assertion en sens inverse résulte de [$\qquad$, 1.2.]. Pour (ii), on sait déjà que lin. équiv. à $0\Longrightarrow$ alb-équiv. à $0$, prouvons le en sens inverse. On doit prouver que si

% archive sheet number: 28
\scanpage{30}
$Z$ est un diviseur sur $X\times T$ et $\unsure{\mathfrak{a}|\alpha|a|\mathfrak{t}}$ sur $T$ est de degré $0$ et tel que $S_T(\mathfrak{a})=0$, alors $Z(\mathfrak{a})$ est lin. équiv. à $0$. Or $\unsure{c^1|c'|\text{on a}}Z(\mathfrak{a})=u(S_T(\mathfrak{a}))$, où $u:\deleted{\mathrm{Alb}}\;A_T\longrightarrow\underline{\mathrm{Pic}}_X^0$ est l'hom défini par $Z$.

\uline{Théorème} 1.\deleted{\unsure{5}}6. Soient $z,z'$ deux classes dans $A'{}^*(X)$ alg. équiv. à $0$, $X$ étant \deleted{[projective]}\inserted{qu. pr.}. Alors $zz'=0$. [Si $X$ est projective, alors il suffit que l'un des \deleted{facteurs} $z$ soit alg. équiv. à $0$, et l'autre \uline{num.} équiv. à $0$ ??]

% slanted
\margin{à voir \unsure{si il|s'il} ne \unsure{suffirait|suffit} pas que $z$ ou $z'$ ne soit que \unsure{$\tau$-équiv.|$z$-équiv.|num. équiv.} à $0$, voire (si $X$ \unsure{projectif|projec.}) équiv. \uline{num.} à $0$.}

\deleted{Cela}\inserted{La première assertion} résulte formellement de la définition de ``alg. équiv. à $0$'' \inserted{via variétés $T$ \uline{projectives} \uline{lisses} connexes,} et du fait que si $T,T'$ sont deux \inserted{telles} variétés \deleted{et} (lisse connexe suffit) et $x,y$ sont deux points de $T$, $x'$ et $y'$ deux pts de $T'$, alors $\mathfrak{a}=\inserted{((x)-(y))\times((x')-(y'))=}(x,x')-(x,y')-(y,x')+(y,y')$ est tel que $\deleted{\unsure{S_T(\mathfrak{a})|S_{\mathfrak{a}}}}\;S_{T\times T'}(\mathfrak{a})=0$.

\margin{des courbes si on veut !}

% slanted
\margin{NB. On ignore si cet $\mathfrak{a}$ est lin. équiv. à $0$, $\overline{\mathrm{m}}$ si $T,T'$ sont des \uline{courbes} !}

[Il y a intérêt à transformer le \deleted{Pb}\inserted{l'énoncé} en un énoncé sur l'opération $\times$]

% archive sheet number: 29
